\begin{figure}[ht]
		\centering
		\begin{tikzpicture}[>=Stealth, line width=0.8pt]
			
			% --- coordinate dei nodi ---
			\coordinate (A) at (0,0);
			\coordinate (B) at (3,0);
			\coordinate (C) at (0,2);
			\coordinate (D) at (1.5,2);
			\coordinate (E) at (3,2);
			\coordinate (F) at (5,2);
			\coordinate (G) at (0,4);
			\coordinate (H) at (3,4);
			
			% --- tratti di trave (linee continue, attraversate ai nodi rigidi C ed E) ---
			\draw[line width=1.5pt] (A) -- (G);   % colonna sinistra A-C-G
			\draw[line width=1.5pt] (B) -- (H);   % colonna destra B-E-H
			\draw[line width=1.5pt] (C) -- (F);   % traversa intermedia C-D-E-F
			\draw[line width=1.5pt] (G) -- (H);   % traversa superiore G-H
			%\draw[line width=1.5pt] (D) -- (H);   % diagonale D-H
			% --- diagonale D-H, con cerniera all'estremità lato H ---
			\coordinate (Hd) at ($(H)!0.056!(D)$);
			%\coordinate (Hd) at ($(H)!0.10!(D)$);   % punto lungo DH, vicino a H
			\draw[line width=1.5pt] (D) -- (Hd);    % asta DH fino al cerchio
			\draw[fill=white, thick] (Hd) circle (2.5pt);   % perno della cerniera
			
			% --- cerniera al suolo in A ---
			\draw[thick] (A) -- ($(A)+(-0.22,-0.38)$)
			-- ($(A)+(0.22,-0.38)$) -- cycle;
			\draw[thick] ($(A)+(-0.32,-0.38)$) -- ($(A)+(0.32,-0.38)$);
			\foreach \x in {-0.24,-0.12,0,0.12,0.24}{
				\draw[thin] ($(A)+(\x,-0.38)$) -- ($(A)+(\x-0.10,-0.52)$);
			}
			\node[circle, fill=white, draw=black, inner sep=0pt,
			minimum size=5pt, line width=0.8pt] at (A) {};
			
			% --- cerniera al suolo in B ---
			\draw[thick] (B) -- ($(B)+(-0.22,-0.38)$)
			-- ($(B)+(0.22,-0.38)$) -- cycle;
			\draw[thick] ($(B)+(-0.32,-0.38)$) -- ($(B)+(0.32,-0.38)$);
			\foreach \x in {-0.24,-0.12,0,0.12,0.24}{
				\draw[thin] ($(B)+(\x,-0.38)$) -- ($(B)+(\x-0.10,-0.52)$);
			}
			\node[circle, fill=white, draw=black, inner sep=0pt,
			minimum size=5pt, line width=0.8pt] at (B) {};
			
			% --- carrello in F (asse verticale, vincolo sotto la trave) ---
			\draw[thick] (F) -- ($(F)+(-0.22,-0.38)$)
			-- ($(F)+(0.22,-0.38)$) -- cycle;
			\draw[thick] ($(F)+(-0.15,-0.44)$) circle (0.07);
			\draw[thick] ($(F)+(0.15,-0.44)$) circle (0.07);
			\draw[thick] ($(F)+(-0.32,-0.52)$) -- ($(F)+(0.32,-0.52)$);
			\foreach \x in {-0.24,-0.12,0,0.12,0.24}{
				\draw[thin] ($(F)+(\x,-0.52)$) -- ($(F)+(\x-0.10,-0.66)$);
			}
			\node[circle, fill=white, draw=black, inner sep=0pt,
			minimum size=5pt, line width=0.8pt] at (F) {};
			
			% --- cerniere interne (cerchi bianchi) in D, G, H ---
			\draw[fill=white, thick] (D) circle (2.5pt);
			\draw[fill=white, thick] (G) circle (2.5pt);
			%\draw[fill=white, thick] (H) circle (2.5pt);
			
			% --- etichette dei nodi ---
			\node[font=\scriptsize] at ($(A)+(-0.22,0)$) {$A$};
			\node[font=\scriptsize] at ($(B)+(0.22,0)$) {$B$};
			\node[font=\scriptsize] at ($(C)+(-0.20,0)$)  {$C$};
			\node[font=\scriptsize] at ($(D)+(0,-0.22)$)  {$D$};
			\node[font=\scriptsize] at ($(E)+(0.22,-0.18)$) {$E$};
			\node[font=\scriptsize] at ($(F)+(0,0.22)$)   {$F$};
			\node[font=\scriptsize] at ($(G)+(-0.20,0.18)$) {$G$};
			\node[font=\scriptsize] at ($(H)+(0.22,0.18)$)  {$H$};
			
		\end{tikzpicture}
		\caption{Esempio di struttura a connessioni eterogenee.
			Vincoli esterni: cerniere al suolo in $A$ e $B$,
			carrello verticale in $F$. La traversa intermedia
			$C$--$D$--$E$--$F$ è continua attraverso $C$ ed $E$
			ed è sconnessa in $D$ con cerniera interna. La
			traversa $G$--$H$ termina in $G$ con cerniera
			d'estremità. In $H$ la traversa $G$--$H$ e il tratto
			$E$--$H$ sono continui, mentre la diagonale $D$--$H$
			vi si connette con cerniera d'estremità.}
		\label{fig:struttura_mista}
	\end{figure}
